\section{Dual Antiplatelet Therapy after Percutaneous Coronary Intervention}

Percutaneous Coronary Intervention (PCI) is a widely used procedure for
treating coronary artery disease. It involves the mechanical widening of
a narrowed coronary artery, typically followed by the deployment of a
drug-eluting stent (DES) to maintain vessel patency. After stent
implantation, patients are at risk of stent thrombosis — a potentially
life-threatening acute occlusion — that is mitigated by Dual Antiplatelet
Therapy (DAPT), combining aspirin with a P2Y12 inhibitor (such as clopidogrel,
ticagrelor, or prasugrel).

The optimal duration of DAPT after PCI is one of the most debated topics in
interventional cardiology. Longer DAPT reduces ischemic risk (particularly
stent thrombosis and myocardial infarction) but simultaneously increases the
risk of bleeding complications. Major bleeding events are not merely
nuisance side effects: they are independently associated with increased
mortality, rehospitalization, and quality-of-life impairment. The trade-off
between ischemic protection and bleeding harm is therefore central to the
clinical decision.

\section{Why Not Every Patient Benefits Equally}

The average-level conclusion of MASTER DAPT — that abbreviated therapy is
noninferior for ischemic events and superior for bleeding — masks a clinically
plausible and important hypothesis: that the optimal strategy differs across
patients. Consider two contrasting patient archetypes:
\begin{itemize}
  \item A patient aged 85 with severe chronic kidney disease, active oral
    anticoagulation, and a recent gastrointestinal bleed: the dominant risk
    is bleeding, and abbreviated DAPT may be strongly beneficial.
  \item A 68-year-old diabetic patient who recently had a myocardial
    infarction, with a complex bifurcation stenting procedure: the dominant
    risk is ischemia, and prolonged DAPT may be preferable.
\end{itemize}

These archetypes illustrate treatment effect heterogeneity that cannot be
captured by ATE analysis but is the target of CATE estimation. Moreover, the
HBR population of MASTER DAPT is particularly heterogeneous: the mean number
of HBR criteria was $2.1 \pm 1.1$ per patient, spanning elderly frail patients,
patients on anticoagulation, and patients with multiple comorbidities.
Whether this hypothesis holds in practice is tested formally in Chapters 6–7.

\todo{Confirm which baseline covariates are available in the MASTER DAPT
dataset for CATE estimation.}
\todo{Confirm whether the thesis uses the real MASTER DAPT individual patient
dataset, a simulated surrogate, or a pre-processed anonymized version.}